% According to ISO, holes should be at a distance of 1.2cm from the border. Hence, an inner border of 2*1.2cm
\usepackage[margin=1.6cm,includefoot,inner=2.4cm]{geometry}

\usepackage[utf8]{inputenc}

\usepackage{helvet}
\usepackage[cm]{sfmath}
\renewcommand{\familydefault}{\sfdefault}

\usepackage{enumitem}

\usepackage{graphicx}

\usepackage{eso-pic}
\usepackage{tikz}

% Model the background
\AddToShipoutPictureBG{ \ifthispageodd { \put(-6,0){
\begin{tikzpicture}
	\node[minimum height=\paperheight,minimum width = 8.1cm,anchor=west,fill=black!20] at (0,0) {};
\end{tikzpicture}
} } {
\begin{tikzpicture}
	\node at (0,0) {};
	\node[minimum height=\paperheight,minimum width = 8.1cm,anchor=west,fill=black!20] at (13.1cm,0) {};
\end{tikzpicture}
} }

\newcounter{recipe}[section]

% Usage: \begin{recipe}{Leckerer Kuchen}{\begin{ingredientcat}...\end{ingredientcat}}{\step{}...\hints{}}
\newcommand{\recipe}[4][]
	{ \begin{minipage}{\textwidth}
		{\noindent\huge \textbf{#2}} \hfill \stepcounter{recipe}\therecipe \\ \rule{\textwidth}{1pt}
		\def\temp{#1}\ifx\temp\empty \\[\bigskipamount] \else \begin{flushright} \vspace{-\medskipamount} \small Ergibt #1 \end{flushright} \fi
		\ifthispageodd {
			\begin{minipage}[t]{.3\textwidth} #3 \end{minipage} \hspace{.03\textwidth}
			\begin{minipage}[t]{.65\textwidth}
				 \noindent {\large \textbf{Anleitung}} \\[-\medskipamount] \rule{\textwidth}{.5pt} 
				\begin{enumerate} #4 \end{enumerate}
			\end{minipage} \\[\bigskipamount]
		} {
			\begin{minipage}[t]{.65\textwidth}
				\noindent {\large \textbf{Anleitung}} \\[-\medskipamount] \rule{\textwidth}{.5pt} 
				\begin{enumerate} #4 \end{enumerate}
			\end{minipage}
			\hspace{.03\textwidth}
			\begin{minipage}[t]{.3\textwidth} #3 \end{minipage} \\[\bigskipamount]
		}
	\end{minipage} }
	
	
% Usage: \yields{1 Kuchen}
% Must be used directly after \begin{recipe}{...}
\newcommand{\yields}[1]
	{\begin{flushright} 
		\vspace{-\medskipamount} \small Ergibt #1
	\end{flushright} }

% Usage: \begin{category}{Teig} \ingredient... \end{category}
\newenvironment{category}[1][]
	{ \def\temp{#1}\ifx\temp\empty \else {\large \textbf{#1}} \fi \begin{itemize}[noitemsep] }
 	{ \end{itemize} }
	
% Usage: \hints{Der Ofen wird angemacht, indem man den Pröngel dreht}
\newcommand{\hints}
	{ \end{enumerate} \noindent {\large \textbf{Hinweise}} \\[-\medskipamount] \rule{\textwidth}{.5pt} \begin{enumerate} \item[]
	}
	
% Usage: \step{Backen}{Backofen öffnen und Zeug reinschieben}
\newcommand{\step}[2][] { \item \def\temp{#1}\ifx\temp\empty \else \textbf{#1} \fi #2 }

% Usage: \ingredient{200}{g}{Zucker} 	
\newcommand{\ingredient}[3]{%
	\item%
	\def\temp{#1}\ifx\temp\empty \else $#1$ \fi 
	\def\temp{#2}\ifx\temp\empty \else #2 \fi
	#3
}